% tikz-matrixmath-doc.tex
% arara: lualatex: { options: ["-shell-escape"] }
% arara: lualatex: { options: ["-shell-escape"] }
\PassOptionsToPackage{english}{babel}
\PassOptionsToPackage{colorlinks = true}{hyperref}
\PassOptionsToPackage{margin = 1in}{geometry}
\PassOptionsToPackage{depth = 4}{bookmark}
\documentclass[12pt]{report}
\usepackage{
    tikz-matrixmath
    ,microtype
    ,babel
    ,fontspec
    ,unicode-math,
    ,caption
    ,subcaption
    ,xcolor
    ,geometry
    ,minted
    ,bookmark
    ,hyperref
}
\setmathfont{IBMPlexMath-Regular}
\counterwithin{figure}{section}
\counterwithin{table}{section}
\DeclareCaptionFormat{tikz-matrixmath}{#1#2\par\noindent#3}
\captionsetup{
    justification = raggedright
    ,singlelinecheck = false
    ,format = tikz-matrixmath
}
\captionsetup[subfigure]{format=hang}
\setcounter{tocdepth}{3}
\setmainfont{Carlito}
\title{
    The tikz-matrixmath Package,\\
    Version 1.0.0\\
    {\small https://github.com/Pseudonym321/TikZ-Animations/tree/master1/TikZ/tikz-matrixmath}
}
\author{Jasper}
\date{\today}
\begin{document}
    \maketitle 
    \tableofcontents
    \chapter{Introduction}
        tikz-matrixmath is a Lua-based 3D graphics package for 
        Ti\textit{k}Z.
        \section{What are the current capabilities of tikz-matrixmath?}
            tikz-matrixmath can render multiple non-intersecting 
            triangulated mesh surfaces at once. It is also possible 
            to draw a plane which has been clipped by a rectangular
            prism.
            There are commands for drawing and labelling geometric 
            vectors too. Curves and surfaces can be defined using 
            matrix transformations.
        \section{
            How does tikz-matrixmath improve on existing packages?
        }
            There are two significant packages for 3D artwork in 
            Ti\textit{k}Z: pgfplots and tikz-3dplot.
            pgfplots can z-sort the faces of a single parametric 
            surface, but not multiple surfaces at once.
            Both pgfplots and tikz-3dplot are capable of changing 
            the camera view, but the result is necessarily 
            orthographic, and the rotations are restricted to 
            azimuth and elevation; that is, the \(z\)-axis is always 
            vertical. tikz-3dplot offers useful pgfmath capabilities 
            for doing vector operations.
            Neither pgfplots nor tikz-3dplot enable 
            matrix transformations. tikz-matrixmath is built to render 
            multiple parametric objects at once, enable arbitrary
            object orientations in non-orthographic projections, 
            and enable matrix transformations.
        \section{Where is the package at, and where it is heading?}
            The vision of the package tikz-matrixmath is to one 
            day be able to render arbitrary numbers of intersecting 
            polygons, by clipping them with each other; this would 
            enable arbitrary surface intersections, without 
            visual artefacts. A more near term goal of 
            tikz-matrixmath is to also implement an ODE 
            solving tool which approximates trajectories in 
            vector fields. I also want to add capabilities for 
            doing programming in Lua, instead of in pgf,
            essentially replacing the pgfmath capabilities of 
            tikz-3dplot with faster Lua-based capabilities.
            Soon there will also be proper camera culling for 
            projective scenes; that is, the action of not drawing 
            segments which are behind the viewing plane.\par 
            There iss no way  to store macros in Lua from the \TeX\      
            end yet, so sometimes there is a problem with 
            repetition. While the geometric vectors are nice,
            they are not yet rendered in 3D like the parametric 
            objects---I want to fix this as well.
            I will also implement a way for the user to create 
            their own reusable Lua functions.
    \chapter{Geometric Vectors}
        \section{How are geometric vectors defined in tikz-matrixmath?}
            Geometric vectors are defined by two things: a 
            start-coordinate, and a relative displacement.
            There are two classifications of geometric vectors;
            there are zero vectors, and non-zero vectors.
            A zero vector is represented by a small circle.
            A non-zero vector is represented by a line segment
            with an arrow. The endpoints of a non-zerovector 
            may or may not have a circle centered on the endpoint.
            If there is no circle drawn, the arrow stops directly
            at the endpoint. If a circle is drawn, then the arrow 
            tip is moved back by the circle's radius. Observe 
            the tips of the vectors in Figure \ref{fig:2-1-1}.
            Both tips lie on the same point, despite one arrow tip 
            being offset by the circle's radius.
            \begin{figure}
                \centering
                \begin{tikzpicture}
                    \drawvector[points = both,green]{0,0}{2,0}
                    \drawvector[blue]{4,0}{-2,0}
                \end{tikzpicture}
                \caption{The exact location of a geometric vector's tip.}
                \label{fig:2-1-1}
            \end{figure}
            The syntax for drawing Figure \ref{fig:2-1-1} is shown in 
            the following code.
\begin{minted}{tex}
\begin{tikzpicture}
    \drawvector[points = both,yellow]{0,0}{2,0}
    \drawvector[blue]{4,0}{-2,0}
\end{tikzpicture}
\end{minted}
            There are four values for the key 
            \mintinline{latex}{points}: \mintinline{latex}{neither}, 
            \mintinline{latex}{behind}, \mintinline{latex}{front} and 
            \mintinline{latex}{both}. These tell the command
            \mintinline{latex}{\drawvector} whether and where to put 
            circles over the endpoints. See Figure  \ref{fig:2-1-2}, 
            which uses  the following code.
\begin{minted}{tex}
\begin{tikzpicture}
    \drawvector[points = neither]{0,0}{0,1} % a)
    \drawvector[points = behind]{1,0}{0,1} % b)
    \drawvector[points = front]{2,0}{0,1} % c)
    \drawvector[points = both]{3,0}{0,1} % d)
\end{tikzpicture}
\end{minted}
            \begin{figure}
                \centering
                \begin{subfigure}{0.3\textwidth}
                    \centering
                    \begin{tikzpicture}
                        \drawvector[points = neither]{0,0}{0,1}
                    \end{tikzpicture}
                    \caption{\mintinline{latex}{points = neither}}
                \end{subfigure}
                \begin{subfigure}{0.3\textwidth}
                    \centering
                    \begin{tikzpicture}
                        \drawvector[points = behind]{1,0}{0,1}
                    \end{tikzpicture}
                    \caption{\mintinline{latex}{points = behind}}
                \end{subfigure}\par
                \begin{subfigure}{0.3\textwidth}
                    \centering
                    \begin{tikzpicture}
                        \drawvector[points = front]{2,0}{0,1}
                    \end{tikzpicture}
                    \caption{\mintinline{latex}{points = front}}
                \end{subfigure}
                \begin{subfigure}{0.3\textwidth}
                    \centering
                    \begin{tikzpicture}
                        \drawvector[points = both]{3,0}{0,1}
                    \end{tikzpicture}
                    \caption{\mintinline{latex}{points = both}}
                \end{subfigure}
                \caption{The values for the key \mintinline{latex}{points}.}
                \label{fig:2-1-2}
            \end{figure}
            The \mintinline{latex}{\drawvector} syntax is mostly
            meant for diagrams which don't use \(z\)-sorting.
            Still, they can be used in some 3D still-diagrams;
            for instance, these vectors connect nicely.
            They will also eventually be able to be drawn flat 
            on planes. There is also a zero-vector command 
            available via \mintinline{latex}{\drawpoint}.
            See Figure \ref{fig:2-1-3}, with the following code.
\begin{minted}{tex}
\begin{tikzpicture}
    \drawpoint[red]{0,0}
\end{tikzpicture}
\end{minted}
            \begin{figure}
                \centering
                \begin{tikzpicture}
                    \drawpoint[red]{0,0}
                \end{tikzpicture}
                \caption{Zero-vector}
                \label{fig:2-1-3}
            \end{figure}
 
    \chapter{The Math Module}
        When the user appends a parametric object to the segments 
        list, they must specify several attributes for the object.
        In particular, when they specify the parametric equations 
        for the object, this is done using Lua syntax.
        Specifically, the user writes Lua code into a tikz key, 
        and that Lua code is later \mintinline{lua}{load()}ed by 
        Lua.  A useful feature of the lua function 
        \mintinline{lua}{load()} is that it allows the user to 
        input functions without their usual table prefix;
        this means that the user can write just 
        \mintinline{lua}{cos()}, instead of 
        \mintinline{lua}{math.cos()}. tikz-matrixmath defines a growing
        handful of math operations, from elementary linear matrix 
        transformations to stereographic projections.
        This list could always be expanded though as needs arise,
        so to make it robust, I will make a way for the user to 
        make their own. See Chapter 1.\par 
        \textbf{tikz-matrixmath uses row-vector convention!} This means that 
        you put the transformation on the right of the matrix 
        multiplication unlike with column vectors. Additionally,
        \textbf{tikz-matrixmath uses matrix-vector convention!} This means 
        that vectors are formatted like matrices---meaning that the
        vector table is nested within a matrix table. This makes it 
        easier for the tools to communicate, and expains why we 
        need to index our functions at depth 2 instead of depth 1.
        \section{The math functions.}
            In addition to those provided by the math module, 
            these ones are also added.
            \begin{itemize}
                \item \mintinline{lua}{matrix_multiply(A, B)}
                \item \mintinline{lua}{matrix_scale(factor, A)}
                \item \mintinline{lua}{reciprocate_by_homogenous(vector)}
                \item \mintinline{lua}{matrix_add(A, B)}
                \item \mintinline{lua}{matrix_subtract(A, B)}
                \item \mintinline{lua}{transpose(A)}
                \item \mintinline{lua}{inverse(matrix)}
                \item \mintinline{lua}{det(matrix)}
                \item \mintinline{lua}{yrotation(angle)}
                \item \mintinline{lua}{translate(x, y, z)}
                \item \mintinline{lua}{xscale(scale)}
                \item \mintinline{lua}{yscale(scale)}
                \item \mintinline{lua}{zscale(scale)}
                \item \mintinline{lua}{scale(scale)}
                \item \mintinline{lua}{xrotation(angle)}
                \item \mintinline{lua}{zrotation(angle)}
                \item \mintinline{lua}{euler(alpha, beta, gamma)}
                \item \mintinline{lua}{sphere(longitude, latitude)}
                \item \mintinline{lua}{dot_product(u, v)}
                \item \mintinline{lua}{cross_product(u, v)}
                \item \mintinline{lua}{norm(u)}
                \item \mintinline{lua}{normalize(u)}
                \item \mintinline{lua}{identity_matrix()}
                \item \mintinline{lua}{stereographic_projection(point)}
            \end{itemize}
    \chapter{Parametric Objects}
        \section{How is a parametric object defined?}
            Parametric objects are defined using a handful of 
            different pieces of information. In particular, they 
            use parametric functions, domain parameters and a samples
            parameter for each dimension in the domain. Of course, 
            there are also parameters for color and such as well. 
            Figure \ref{fig:3-1-1} is defined by the following code:
\begin{minted}{tex}
\begin{tikzpicture}
    \appendcurve[
        u min = 0
        ,u max = {tau}
        ,u samples = 200
        ,transformation = {euler(pi/2,pi/3,pi/2)}
        ,x = {2*cos(u)}
        ,y = {2*sin(u)}
        ,z = {u}
        ,draw options = {
            purple
            ,ultra thick
            ,line cap = round
        }
    ]
    \appendsurface[
        u min = 0
        ,u max = {tau}
        ,u samples = 36
        ,v min = 0
        ,v max = {tau}
        ,v samples = 36*2/3
        ,transformation = {euler(pi/2,pi/3,pi/2)}
        ,x = {2*cos(u)+sphere(u,v)[1][1]}
        ,y = {2*sin(u)+sphere(u,v)[1][2]}
        ,z = {u+sphere(u,v)[1][3]}
        ,draw options = {
            draw
            ,blue,ultra thin
            ,line join = round
            ,line cap = round
        }
        ,fill options = {
            fill = green
            ,fill opacity = 0.7
        }
    ]
    \foreach \u in {0,...,35} {
        \foreach \v in {0,...,35} {
            \pgfmathsetmacro{\uscale}{2*pi/(36-1)}
            \pgfmathsetmacro{\vscale}{2*pi/(36*2/3-1)}
            \pgfmathsetmacro{\u}{\u*\uscale}
            \pgfmathsetmacro{\v}{\v*\vscale}
            \appendcurve[
                u min = 0
                ,u max = 0.3
                ,u samples = 2
                ,transformation = {euler(pi/2,pi/3,pi/2)}
                ,x = {
                    2 * 
                    cos(token.get_macro("u")) +
                    (u + 1) * 
                    sphere(
                        token.get_macro("u")
                        ,token.get_macro("v")
                    )[1][1]
                }
                ,y = {
                    2 * 
                    sin(token.get_macro("u")) +
                    (u + 1) * 
                    sphere(
                        token.get_macro("u")
                        ,token.get_macro("v")
                    )[1][2]
                }
                ,z = {
                    token.get_macro("u") + 
                    (u + 1) * 
                    sphere(
                        token.get_macro("u")
                        ,token.get_macro("v")
                    )[1][3]
                }
                ,draw options = {
                    -{Stealth[round]}
                    ,line cap = round
                }
            ]
        } 
    }
    \rendersegments
\end{tikzpicture}
\end{minted}
            \begin{figure}
                \centering
                \begin{tikzpicture}
                    \appendcurve[
                        u min = 0
                        ,u max = {tau}
                        ,u samples = 200
                        ,transformation = {euler(pi/2,pi/3,pi/2)}
                        ,x = {2*cos(u)}
                        ,y = {2*sin(u)}
                        ,z = {u}
                        ,draw options = {
                            purple
                            ,ultra thick
                            ,line cap = round
                        }
                    ]
                    \appendsurface[
                        u min = 0
                        ,u max = {tau}
                        ,u samples = 36
                        ,v min = 0
                        ,v max = {tau}
                        ,v samples = 36*2/3
                        ,transformation = {euler(pi/2,pi/3,pi/2)}
                        ,x = {2*cos(u)+sphere(u,v)[1][1]}
                        ,y = {2*sin(u)+sphere(u,v)[1][2]}
                        ,z = {u+sphere(u,v)[1][3]}
                        ,draw options = {
                            draw
                            ,blue,ultra thin
                            ,line join = round
                            ,line cap = round
                        }
                        ,fill options = {
                            fill = green
                            ,fill opacity = 0.7
                        }
                    ]
                    \foreach \u in {0,...,35} {
                        \foreach \v in {0,...,35} {
                            \pgfmathsetmacro{\uscale}{2*pi/(36-1)}
                            \pgfmathsetmacro{\vscale}{2*pi/(36*2/3-1)}
                            \pgfmathsetmacro{\u}{\u*\uscale}
                            \pgfmathsetmacro{\v}{\v*\vscale}
                            \appendcurve[
                                u min = 0
                                ,u max = 0.3
                                ,u samples = 2
                                ,transformation = 
                                    {euler(pi/2,pi/3,pi/2)}
                                ,x = {
                                    2 * 
                                    cos(token.get_macro("u")) +
                                    (u + 1) * 
                                    sphere(
                                        token.get_macro("u")
                                        ,token.get_macro("v")
                                    )[1][1]
                                }
                                ,y = {
                                    2 * 
                                    sin(token.get_macro("u")) +
                                    (u + 1) * 
                                    sphere(
                                        token.get_macro("u")
                                        ,token.get_macro("v")
                                    )[1][2]
                                }
                                ,z = {
                                    token.get_macro("u") + 
                                    (u + 1) * 
                                    sphere(
                                        token.get_macro("u")
                                        ,token.get_macro("v")
                                    )[1][3]
                                }
                                ,draw options = {
                                    -{Stealth[round]}
                                    ,line cap = round
                                }
                            ]
                        } 
                    }
                    \rendersegments
                \end{tikzpicture}
                \caption{A triangulated mesh surface and a curve.}
                \label{fig:3-1-1}
            \end{figure}
            tikz-matrixmath is also capable of handling division by zero
            without erroring. This is because the math is handled in
            Lua, where division by zero doesn't error. We clip our 
            results to be well within the scope of the Ti\textit{z}Z
            canvas. The code for Figure \ref{fig:3-1-2} is as follows.
\begin{minted}{tex}
\begin{tikzpicture}
    \clip (-5,-5) rectangle (5,5);
    \pgfmathsetmacro{\angle}{0}
    \let\angle\angle
    \foreach \longitude in {1,...,36}{
        \let\longitude\longitude
        \appendcurve[
            u min = 0
            ,u max = {pi}
            ,u samples = 4*45
            ,transformation = {euler(pi/2,pi/3,pi/6)}
            ,x = {
                stereographic_projection(
                    matrix_multiply(
                        sphere(
                            tau * 
                            token.get_macro("longitude") /
                            36
                            ,u
                        )
                        ,euler(
                            0
                            ,pi/2
                            ,(2*pi/24) *
                            token.get_macro("angle") / 
                            36
                        )
                    )
                )[1][1]
            }
            ,y = {
                stereographic_projection(
                    matrix_multiply(
                        sphere(
                            tau *
                            token.get_macro("longitude") /
                            36
                            ,u
                        )
                        ,euler(
                            0
                            ,pi/2
                            ,(2*pi/24) * 
                            token.get_macro("angle") / 
                            36
                        )
                    )
                )[1][2]
            }
            ,z = 0*u
        ]
        \appendcurve[
            u min = 0
            ,u max = {pi}
            ,u samples = 23
            ,transformation = {
                matrix_multiply(
                    euler(
                        0
                        ,pi/2
                        ,(2*pi/24) * 
                        token.get_macro("angle") / 
                        36
                    )
                    ,euler(pi/2,pi/3,pi/6)
                )
            }
            ,x = {
                sphere(
                    tau * 
                    token.get_macro("longitude") / 
                    36
                    ,u
                )[1][1]
            }
            ,y = {
                sphere(
                    tau * 
                    token.get_macro("longitude") / 
                    36
                    ,u
                )[1][2]
            }
            ,z = {
                sphere(
                    tau * 
                    token.get_macro("longitude") / 
                    36
                    ,u
                )[1][3]
            }
        ]
    }
    \foreach \latitude in {1,...,18}{
        \let\latitude\latitude
        \appendcurve[
            u min = 0
            ,u max = {2*pi}
            ,u samples = 4*45
            ,transformation = {euler(pi/2,pi/3,pi/6)}
            ,x = {
                stereographic_projection(
                    matrix_multiply(
                        sphere(
                            u
                            ,tau *
                            token.get_macro("latitude") / 
                            36
                        )
                        ,euler(
                            0
                            ,pi/2
                            ,(tau/24) * 
                            token.get_macro("angle") / 
                            36
                        )
                    )
                )[1][1]
            }
            ,y = {
                stereographic_projection(
                    matrix_multiply(
                        sphere(
                            u
                            ,tau *
                            token.get_macro("latitude") /
                            36
                        )
                        ,euler(
                            0
                            ,pi/2
                            ,(tau/24) *
                            token.get_macro("angle") /
                            36
                        )
                    )
                )[1][2]
            }
            ,z = 0*u
        ]
        \appendcurve[
            u min = 0
            ,u max = {2*pi}
            ,u samples = 45
            ,transformation = {
                matrix_multiply(
                    euler(
                        0
                        ,pi/2
                        ,(2*pi/24) *
                        token.get_macro("angle") /
                        36
                    )
                    ,euler(pi/2,pi/3,pi/6)
                )
            }
            ,x = {
                sphere(
                    u
                    ,tau*token.get_macro("latitude")/36
                )[1][1]
            }
            ,y = {
                sphere(
                    u
                    ,tau*token.get_macro("latitude")/36
                )[1][2]
            }
            ,z = {
                sphere(
                    u
                    ,tau*token.get_macro("latitude")/36
                )[1][3]
            }
        ]
    }
    \rendersegments
\end{tikzpicture}
\end{minted}
            \begin{figure}
                \centering
                \begin{tikzpicture}
                    \clip 
                        (-\textwidth/2,-5) 
                        rectangle 
                        (\textwidth/2,5)
                    ;
                    \pgfmathsetmacro{\angle}{0}
                    \let\angle\angle
                    \foreach \longitude in {1,...,36}{
                        \let\longitude\longitude
                        \appendcurve[
                            u min = 0
                            ,u max = {pi}
                            ,u samples = 4*45
                            ,transformation = {euler(pi/2,pi/3,pi/6)}
                            ,draw options = {
                                line cap = round
                            }
                            ,x = {
                                stereographic_projection(
                                    matrix_multiply(
                                        sphere(
                                            tau * 
                                            token.get_macro("longitude") /
                                            36
                                            ,u
                                        )
                                        ,euler(
                                            0
                                            ,pi/2
                                            ,(2*pi/24) *
                                            token.get_macro("angle") / 
                                            36
                                        )
                                    )
                                )[1][1]
                            }
                            ,y = {
                                stereographic_projection(
                                    matrix_multiply(
                                        sphere(
                                            tau *
                                            token.get_macro("longitude") /
                                            36
                                            ,u
                                        )
                                        ,euler(
                                            0
                                            ,pi/2
                                            ,(2*pi/24) * 
                                            token.get_macro("angle") / 
                                            36
                                        )
                                    )
                                )[1][2]
                            }
                            ,z = 0*u
                        ]
                        \appendcurve[
                            u min = 0
                            ,u max = {pi}
                            ,u samples = 23
                            ,transformation = {
                                matrix_multiply(
                                    euler(
                                        0
                                        ,pi/2
                                        ,(2*pi/24) * 
                                        token.get_macro("angle") / 
                                        36
                                    )
                                    ,euler(pi/2,pi/3,pi/6)
                                )
                            }
                            ,draw options = {
                                line cap = round
                            }
                            ,x = {
                                sphere(
                                    tau * 
                                    token.get_macro("longitude") / 
                                    36
                                    ,u
                                )[1][1]
                            }
                            ,y = {
                                sphere(
                                    tau * 
                                    token.get_macro("longitude") / 
                                    36
                                    ,u
                                )[1][2]
                            }
                            ,z = {
                                sphere(
                                    tau * 
                                    token.get_macro("longitude") / 
                                    36
                                    ,u
                                )[1][3]
                            }
                        ]
                    }
                    \foreach \latitude in {1,...,18}{
                        \let\latitude\latitude
                        \appendcurve[
                            u min = 0
                            ,u max = {2*pi}
                            ,u samples = 4*45
                            ,transformation = {euler(pi/2,pi/3,pi/6)}
                            ,draw options = {
                                line cap = round
                            }
                            ,x = {
                                stereographic_projection(
                                    matrix_multiply(
                                        sphere(
                                            u
                                            ,tau *
                                            token.get_macro("latitude") / 
                                            36
                                        )
                                        ,euler(
                                            0
                                            ,pi/2
                                            ,(tau/24) * 
                                            token.get_macro("angle") / 
                                            36
                                        )
                                    )
                                )[1][1]
                            }
                            ,y = {
                                stereographic_projection(
                                    matrix_multiply(
                                        sphere(
                                            u
                                            ,tau *
                                            token.get_macro("latitude") /
                                            36
                                        )
                                        ,euler(
                                            0
                                            ,pi/2
                                            ,(tau/24) *
                                            token.get_macro("angle") /
                                            36
                                        )
                                    )
                                )[1][2]
                            }
                            ,z = 0*u
                        ]
                        \appendcurve[
                            u min = 0
                            ,u max = {2*pi}
                            ,u samples = 45
                            ,transformation = {
                                matrix_multiply(
                                    euler(
                                        0
                                        ,pi/2
                                        ,(2*pi/24) *
                                        token.get_macro("angle") /
                                        36
                                    )
                                    ,euler(pi/2,pi/3,pi/6)
                                )
                            }
                            ,draw options = {
                                line cap = round
                            }
                            ,x = {
                                sphere(
                                    u
                                    ,tau*token.get_macro("latitude")/36
                                )[1][1]
                            }
                            ,y = {
                                sphere(
                                    u
                                    ,tau*token.get_macro("latitude")/36
                                )[1][2]
                            }
                            ,z = {
                                sphere(
                                    u
                                    ,tau*token.get_macro("latitude")/36
                                )[1][3]
                            }
                        ]
                    }
                    \rendersegments
                \end{tikzpicture}
                \caption{Automatic division by zero handling.}
                \label{fig:3-1-2}
            \end{figure}
            tikz-matrixmath is meant to be compatible with Ti\textit{k}Z, and
            they are meant to be used together---at least for now.
            Maybe one day tikz-matrixmath will have enough capabilities to 
            do things like for loops and conditionals---\textit{at 
            home}. The code for Figure \ref{fig:3-1-3}, which 
            illustrates the use of a conditional, is given.
\begin{minted}{tex}
\begin{tikzpicture}
    \clip 
        (-\textwidth/2,-5) 
        rectangle 
        (\textwidth/2,5)
    ;
    \foreach[count = \c from 1] \angle in {20,40,...,80} {
        \pgfmathsetmacro{\angle}{\angle*pi/180}
        \let\angle\angle
        \appendsurface[
            u min = pi/6
            ,u max = 0.75*tau
            ,u samples = 36 
            ,v min = 0
            ,v max = tau
            ,v samples = 36
            ,transformation = {euler(pi/2,pi/3,pi/6+pi)}
            ,draw options = {
                draw = black
                ,line cap = round
                ,line join = round
            }
            ,fill options = {
                \ifnum\c=1
                    fill = red
                \fi
                \ifnum\c=2
                    fill = yellow
                \fi
                \ifnum\c=3
                    fill = green
                \fi
                \ifnum\c=4
                    fill = blue
                \fi
                ,fill opacity = 0.6
            }
            ,x = {
                (1 / cos(token.get_macro("angle"))) *
                cos(u) + 
                sqrt(
                    (
                        1 / 
                        cos(token.get_macro("angle"))
                    )^2 - cos(token.get_macro("angle"))
                ) * 
                sphere(u,v)[1][1]
            }
            ,y = {
                (1 / cos(token.get_macro("angle"))) *
                sin(u) + 
                sqrt(
                    (
                        1 / 
                        cos(token.get_macro("angle"))
                    )^2 - cos(token.get_macro("angle"))
                ) * 
                sphere(u,v)[1][2]
            }
            ,z = {
                sqrt(
                    (
                        1 / 
                        cos(token.get_macro("angle"))
                    )^2 - cos(token.get_macro("angle"))
                ) * 
                sphere(u,v)[1][3]
            }
        ]
    }
    \rendersegments
\end{tikzpicture}
\end{minted}
            \begin{figure}
                \centering
                \begin{tikzpicture}
                    \clip 
                        (-\textwidth/2,-5) 
                        rectangle 
                        (\textwidth/2,5)
                    ;
                    \foreach[count = \c from 1] \angle in {20,40,...,80} {
                        \pgfmathsetmacro{\angle}{\angle*pi/180}
                        \let\angle\angle
                        \appendsurface[
                            u min = pi/6
                            ,u max = 0.75*tau
                            ,u samples = 36 
                            ,v min = 0
                            ,v max = tau
                            ,v samples = 36
                            ,transformation = {euler(pi/2,pi/3,pi/6+pi)}
                            ,draw options = {
                                draw = black
                                ,line cap = round
                                ,line join = round
                            }
                            ,fill options = {
                                \ifnum\c=1
                                    fill = red
                                \fi
                                \ifnum\c=2
                                    fill = yellow
                                \fi
                                \ifnum\c=3
                                    fill = green
                                \fi
                                \ifnum\c=4
                                    fill = blue
                                \fi
                                ,fill opacity = 0.6
                            }
                            ,x = {
                                (1 / cos(token.get_macro("angle"))) *
                                cos(u) + 
                                sqrt(
                                    (
                                        1 / 
                                        cos(token.get_macro("angle"))
                                    )^2 - cos(token.get_macro("angle"))
                                ) * 
                                sphere(u,v)[1][1]
                            }
                            ,y = {
                                (1 / cos(token.get_macro("angle"))) *
                                sin(u) + 
                                sqrt(
                                    (
                                        1 / 
                                        cos(token.get_macro("angle"))
                                    )^2 - cos(token.get_macro("angle"))
                                ) * 
                                sphere(u,v)[1][2]
                            }
                            ,z = {
                                sqrt(
                                    (
                                        1 / 
                                        cos(token.get_macro("angle"))
                                    )^2 - cos(token.get_macro("angle"))
                                ) * 
                                sphere(u,v)[1][3]
                            }
                        ]
                    }
                    \rendersegments
                \end{tikzpicture}
                \caption{Stereographically nested tori.}
                \label{fig:3-1-3}
            \end{figure}
    \chapter{Clipped Subspaces}
        A big focus for the future of this package is 
        triangle-triangle clipping, which generalizes to 
        polygon clipping. This dream is twofold; firstly it will 
        enable parametric surfaces to intersect, and secondly 
        it will enable intersecting plane diagrams which haven't 
        been achieved in full generality in tikz yet.
        \section{Clipping individual planes.}
            This is where the technology is at currently. We can 
            clip a single plane by a rectangular prism, and draw them.
            Eventually, when the polygon clipping gets going, this will 
            evolve into a tool for graphing intersecting planes without
            using approximations. The code for Figure \ref{fig:4-1-1} 
            is given.
\begin{minted}{tex}
\begin{tikzpicture}
    \appendprism[%
        a = 1
        ,b = 2
        ,c = 3
        ,d = 4
        ,x min = -1
        ,x max = 1
        ,y min = -1
        ,y max = 1
        ,z min = -1
        ,z max = 1
        ,transformation = {euler(pi/2,pi/3,pi/3)}
    ]
    \appendplane[%
        a = 1
        ,b = 2
        ,c = 3
        ,d = 1
        ,x min = -1
        ,x max = 1
        ,y min = -1
        ,y max = 1
        ,z min = -1
        ,z max = 1
        ,transformation = {euler(pi/2,pi/3,pi/3)}
    ]
    \rendersegments
\end{tikzpicture}
\end{minted}
            \begin{figure}
                \centering
                \begin{tikzpicture}
                    \appendprism[%
                        a = 1
                        ,b = 2
                        ,c = 3
                        ,d = 4
                        ,x min = -1
                        ,x max = 1
                        ,y min = -1
                        ,y max = 1
                        ,z min = -1
                        ,z max = 1
                        ,transformation = {euler(pi/2,pi/3,pi/3)}
                    ]
                    \appendplane[%
                        a = 1
                        ,b = 2
                        ,c = 3
                        ,d = 1
                        ,x min = -1
                        ,x max = 1
                        ,y min = -1
                        ,y max = 1
                        ,z min = -1
                        ,z max = 1
                        ,transformation = {euler(pi/2,pi/3,pi/3)}
                    ]
                    \rendersegments
                \end{tikzpicture}
                \caption{%
                    A plane which has been clipped by a rectangular 
                    prism.
                }
                \label{fig:4-1-1}
            \end{figure}
\end{document}